% LaTeX fragment for the original fixed escalation-policy experiment.
% Requires booktabs. No threshold sweep or live-agent results are substituted.
% Source: full_reliability_conformal_escalation rows in
% equi-agent/outputs/fairvision_reliability_selective_arbitration/selective_arbitration_ablation_metrics.csv
% Counts independently checked against selective_arbitration_predictions.csv.

\subsection*{Fixed-policy selective prediction}

We evaluated the original fixed escalation policy of the deterministic
reliability-weighted fusion component on 9,000 Harvard FairVision30K test
cases (3,000 per task). The saved policy designated a case for clinician
review if any of the following conditions held: the fraction of model votes
disagreeing with the majority was at least 0.25; the fused probability lay
within 0.08 of the 0.5 classification threshold; weighted reliability was
below 0.35; or the conformal label set contained other than exactly one label.
The saved conformal configuration used a nominal alpha of 0.10. Cases with
no escalation condition were accepted. This analysis uses the recorded
acceptance decisions and unchanged predictions, rather than a new risk-score
threshold sweep or a rerun of the live agent.

The policy escalated 6,223 of 9,000 cases (69.14\%), retaining 2,777 cases
(30.86\% coverage). Positive-class F1 was 0.7295 when scoring all cases and
0.9020 when scoring only the accepted subset. Within glaucoma, the policy
retained 1,510 cases (50.33\% coverage), with accepted-case F1 of 0.8727.
Within AMD, it retained 1,267 cases (42.23\% coverage), with accepted-case
F1 of 0.9359. Every DR case was escalated, so accepted-case DR F1 was
undefined. Consequently, the pooled accepted-case score contains glaucoma
and AMD cases only; it is not a task-average score or evidence of successful
automated DR classification.

\begin{table*}[t]
\centering
\caption{Performance under the original fixed escalation policy of the
historical deterministic reliability-weighted fusion component on Harvard
FairVision30K. All F1 values are positive-class F1. Full-cohort F1 scores all
saved predictions, including those on cases designated for review;
accepted-case F1 scores only non-escalated cases. Coverage is the percentage
accepted. Pooled F1 is calculated from pooled confusion counts, not averaged
across tasks. All DR cases were escalated, leaving accepted-case DR F1
undefined. These are not results from the live 250-case-per-task RetinAgent
evaluation.}
\label{tab:fairvision_fixed_escalation}
\small
\begin{tabular}{@{}lrrrrcc@{}}
\toprule
Task & Total $n$ & Accepted $n$ & Coverage (\%) & Escalated (\%) &
Full-cohort F1 & Accepted-case F1 \\
\midrule
Glaucoma & 3,000 & 1,510 & 50.33 &  49.67 & 0.7455 & 0.8727 \\
AMD      & 3,000 & 1,267 & 42.23 &  57.77 & 0.7598 & 0.9359 \\
DR       & 3,000 &     0 &  0.00 & 100.00 & 0.5306 & ---    \\
\midrule
Pooled   & 9,000 & 2,777 & 30.86 &  69.14 & 0.7295 & 0.9020 \\
\bottomrule
\end{tabular}
\end{table*}

\paragraph{F1 definition.}
The historical 0.9020 result is positive-class F1, not macro-F1. For
consistency with a macro-F1 manuscript, the same saved confusion counts give
full-cohort and accepted-case macro-F1 values of 0.7568 and 0.8742 for
glaucoma; 0.8022 and 0.9345 for AMD; 0.7398 and undefined for DR; and 0.8004
and 0.9017 when pooling all cases and accepted cases, respectively.
The two metrics should be explicitly distinguished rather than renamed.

\paragraph{Interpretation.}
This demonstrates higher F1 within the subset retained by this particular
historical policy, at substantial loss of coverage. It does not demonstrate
better predictions for the same patients, improved clinical outcomes after
review, or the practicality of referring almost seven in ten cases. The
task and class composition changes between the full and accepted cohorts.
No uncertainty intervals or clinician outcomes are supplied by this
analysis. It also does not establish that language-model reasoning produced
the selective-performance result: the evaluated component is deterministic.
The plot and table should therefore describe a fixed operating point per
task, not a continuous curve or a live-agent experiment. Verifying these
saved results does not independently establish the split purity of every
upstream reliability estimate.
